\subsection{The original tau-base weight line in the exponential reflection flow}\label{endpoint-f1_reflection--the-original-tau-base-weight-line-in-the-exponential-reflection-flow}

13 September 2026. This is a calculation on the existing Split-Zero arithmetic construction, not a replacement finite-field model and not a weight estimate. It joins the original moment line in the tau-base right adjoint to an exact reflection and flow on the already constructed polynomial-exponential family. The finite transport identity extends through the special fibre even though the corresponding contour gauge contains an essential exponential in \(1/u\).

\subsubsection{1. Retained source and conventions}\label{endpoint-f1_reflection--retained-source-and-conventions}

The controlling source is \nolinkurl{Tau_Base_Cohomology_2026-09-12/NOTE.md}, SHA256 \nolinkurl{d03e71ad18f187bd89880cb8ca6fc9c5d6f260b3de8e8e5702c27db97e7b63c3}. Its whole body was read. Sections 3--7 supply, respectively,

\[
\mathsf A_\tau(\mathcal T)
=[V\oplus V\xrightarrow{\Theta(\phi-\widehat\psi)}\mathscr B],
\quad H^1=Q=\mathscr B/\Theta V,
\quad K_\tau(W)\simeq S_\eta W[1].
\tag{RW1}
\]

The original action is \(U_p F(x)=F(x/p)\). Its two source legs are \((U_p\phi,pU_{1/p}\psi)\), not two copies of \(U_p\). Its Mellin multiplier is \(p^s\). The actual second moment \(\int\phi(x)\,dx\) has multiplier \(p\); its k-fold tensor is the specified line \(L_k\) with multiplier \(p^k\). The dual action is consequently \(p^k\) times precomposition by \(U_p^{-1}\). All these maps fix support masks and carry supported zero to supported zero, with external absence tau carried to tau.

Use an actual full quartet packet and its original cyclic tensor module,

\[
E=\mathbb C[S]/(\chi),\quad q=\deg\chi,\quad A=M_S,
\quad S=s_1+\cdots+s_k.
\tag{RW2}
\]

Here the retained quartet and its complete multiplicities give

\[
\overline\chi=\chi,\qquad
\chi(k-S)=\varepsilon\chi(S),\qquad\varepsilon=(-1)^q.
\tag{RW3}
\]

Indeed complex conjugation and reflection permute the tensor-sum roots with their original multiplicities; the leading coefficient of the reflected monic degree-q polynomial is \((-1)^q\). This statement is about those actual symmetry-stable packets. It does not invent an off-line zero or discard a nilpotent direction. The arithmetic inclusion remains \(\eta[P]=\upsilon_h^{\otimes k}P(A_h^{(k)})1\). In particular its full Taylor unit is not absorbed into a scalar or removed from the comparison.

The existing exponential family uses the antiderivative with \(\Phi(0)=0\),

\[
\Phi'=\chi,\quad\Phi_t(S)=\Phi(S)-tS,\quad
D_{u,t}=u\partial_S+\chi(S)-t,
\quad\mathcal H=\operatorname{coker}D_{u,t}.
\tag{RW4}
\]

Leading-term division makes H free of rank q over \(\mathbb C[u,t]\), with basis \(1,S,\ldots,S^{q-1}\). In that basis the already proved connection is

\[
\nabla_t=\partial_t-\frac{A(t)}u,\qquad
A(t)=A+t\mathcal R,
\quad\mathcal R=1\otimes\ell_{q-1}.
\tag{RW5}
\]

The matrix \(A(t)\) is multiplication by S in \(\mathbb C[S]/(\chi-t)\) on these representatives. Multiplication by an arbitrary S-polynomial is not thereby asserted to descend to the full u-deformed differential cokernel. The connection and the chain maps below supply the actual descended operators.

\subsubsection{2. Reflection is a cochain map of the actual family}\label{endpoint-f1_reflection--reflection-is-a-cochain-map-of-the-actual-family}

Let \(RP(S)=P(k-S)\) and

\[
r(u,t)=(-\varepsilon u,\varepsilon t)=(u',t').
\tag{RW6}
\]

The matrix R on degree-below-q representatives has the explicit entries

\[
R_{ij}=\begin{cases}\binom ji k^{j-i}(-1)^i&i\le j,\\0&i>j,
\end{cases}\qquad R^2=I.
\tag{RW7}
\]

For every polynomial P, differentiation and (RW3) give

\[
D_{u',t'}(\varepsilon RP)=R(D_{u,t}P).
\tag{RW8}
\]

In detail, the derivative term on the left is \((-\varepsilon u)\varepsilon(-RP')=uRP'\); its remaining term is \(\varepsilon(\chi-\varepsilon t)RP=(\varepsilon\chi-t)RP\), which is exactly the reflected right side. Thus degree-zero map \(\varepsilon R\) and degree-one map R give a cochain isomorphism, with its inverse given by the same reflection. This extends over \((u,t)=(0,0)\), where it is the original polynomial reflection on E. Adding coefficient conjugation gives the actual dagger operation, with parameter map \((u,t)\mapsto(-\varepsilon\overline u,\varepsilon\overline t)\).

On the companion matrices the same calculation is

\[
R A(t)R=kI-A(\varepsilon t).
\tag{RW9}
\]

This follows either by reducing \((k-S)P(k-S)\) modulo \(\chi-\varepsilon t\), or by using the algebra isomorphism defined by (RW3). All entries and all repeated-root blocks remain; no diagonalization is used.

\subsubsection{3. The weight twist is the original potential's constant, not a choice}\label{endpoint-f1_reflection--the-weight-twist-is-the-original-potentials-constant-not-a-choice}

Integrating (RW3) with the original endpoint gives

\[
\Phi(k-S)=\Phi(k)-\varepsilon\Phi(S),
\] \[
\boxed{\frac{\Phi_t(k-S)}u
=\frac{\Phi_{\varepsilon t}(S)}{-\varepsilon u}
 +\frac{\Phi(k)-kt}{u}.}
\tag{RW10}
\]

Neither \(\Phi(k)\) nor \(k\) is set to zero. Pull the target connection back along r, keeping u fixed while differentiating t. Since \(dt'=\varepsilon\,dt\),

\[
r^*\nabla_t=\partial_t+\frac{A(\varepsilon t)}u,
\quad
(r^*\nabla_t)(Rv)-R\nabla_tv=\frac ku Rv.
\tag{RW11}
\]

Consequently the exact horizontal map on \(u\ne0\) is

\[
\boxed{\mathcal F(u,t)
=\exp\!\left(\frac{\Phi(k)-kt}{u}\right)R.}
\tag{RW12}
\]

Its logarithmic t-derivative is \(-k/u\), so substituting (RW11) proves \((r^*\nabla)(\mathcal Fv)=\mathcal F\nabla v\). This uses the full constant supplied by the original potential. It does not select a new norm on H.

For an oriented rapid-decay contour Gamma, its reflected contour is the pushforward \(\Gamma'=k-\Gamma\) with that orientation. The substitution \(S_{\mathrm{old}}=k-S_{\mathrm{new}}\) has \(dS_{\mathrm{old}}=-dS_{\mathrm{new}}\). Therefore the actual integrals satisfy

\[
\int_\Gamma e^{\Phi_t(S)/u}P(S)\,dS
=-\exp\!\left(\frac{\Phi(k)-kt}{u}\right)
\int_{\Gamma'} e^{\Phi_{t'}(S)/u'}(RP)(S)\,dS.
\tag{RW13}
\]

The minus sign is the one-form pullback. The constant exponential is retained. The exponential identity also carries the original decay along Gamma to Gamma', so this substitution does not postulate a different contour class.

\subsubsection{4. A regular finite flow identity through the special fibre}\label{endpoint-f1_reflection--a-regular-finite-flow-identity-through-the-special-fibre}

The scaling calculation computes the regular lift of forward theta dilation from \(L=u\partial_t-S\). Its coefficient transport C\_a is characterized by

\[
\frac{d}{da}C_a(u,t)=-C_a(u,t)A(t+ua),\qquad C_0=I.
\tag{RW14}
\]

This linear matrix equation, with polynomial coefficients, has a unique entire solution in a, depending holomorphically on u,t. Its local power series recursion is explicit. It obeys

\[
C_{a+b}(u,t)=C_a(u,t)C_b(u,t+ua),
\quad C_a^{-1}(u,t)=C_{-a}(u,t+ua).
\tag{RW15}
\]

To prove the product law, differentiate its right side in b and compare its initial value at b=0 with \(C_{a+b}\) in (RW14). Uniqueness gives equality. In particular this finite transport is invertible even at u=0; there \(C_a(0,t)=\exp(-aA(t))\).

Its reflection is the new exact formula

\[
\boxed{R C_a(u,t)R=e^{-ka}C_{-a}(-\varepsilon u,\varepsilon t).}
\tag{RW16}
\]

Proof: write the left side as B\_a. By (RW9),

\[
B_a'=-kB_a+B_aA(\varepsilon t+\varepsilon ua).
\]

For \(u'=-\varepsilon u\), \(t'=\varepsilon t\), the derivative of \(C_{-a}(u',t')\) is \(C_{-a}(u',t')A(t'-u' a)\). Thus the right side of (RW16) obeys this same equation and has the same initial value I. This proves (RW16) for all parameters, without dividing by u. It is not an attempt to evaluate the singular exponential (RW12) at u=0.

Along \(t\mapsto t+ua\), the ratio of the two full scalar factors in (RW12) is exactly \(\exp(-ka)\). At \(a=\log p\), forward theta dilation has moment line \(p^{-k}\); inverse dilation U\_p has moment line \(p^k\). Thus the very same factor appears both in the actual tau-base adjunction and in the exact exponential-family reflection flow. It is not a finite-field cardinality assigned to the absolute base point.

Specializing (RW16) gives the full arithmetic identity

\[
R p^{-A} R=p^{-k}p^A,
\qquad R p^A R=p^k p^{-A}.
\tag{RW17}
\]

Every block is retained:

\[
p^{\pm A}|_{\lambda}
=p^{\pm\lambda}\sum_{j=0}^{\ell_\lambda-1}
\frac{(\pm\log p)^j}{j!}N_\lambda^j.
\tag{RW18}
\]

\subsubsection{5. The computed adjoint transport and its exact evaluation}\label{endpoint-f1_reflection--the-computed-adjoint-transport-and-its-exact-evaluation}

On coefficient sections the scaling operator is semilinear: \(T_av(t)=C_a(u,t)v(t+ua)\). The k-fold moment line has \(T_a^Lz(t)=\exp(-ka)z(t+ua)\). Therefore the actual twisted dual is

\[
\boxed{T_a^\vee\ell(t)
=e^{-ka}C_a(u,t)^{-T}\ell(t+ua).}
\tag{RW19}
\]

Substitution proves, with no omitted scalar,

\[
(T_a^\vee\ell)(t)^T(T_av)(t)
=e^{-ka}\ell(t+ua)^Tv(t+ua).
\tag{RW20}
\]

The product rule (RW15) proves the composition law for (RW19). Its infinitesimal generator is

\[
u\partial_t+A(t)^T-kI.
\tag{RW21}
\]

At the actual special fibre and inverse dilation it becomes \(p^k(p^A)^{-T}\). This is precisely the action obtained by applying the right adjoint \(K_{(\tau,k)}(L_k)\) in the original source, then precomposing finite jet evaluation. In the conjugate packet, the source's dagger and residue pairing give its specified linear map to this twisted dual. No Hermitian positivity has been inferred from this bilinear calculation.

For completeness, the existing constant residue matrix

\[
\mathsf S_{ij}=[S^{-1}]\frac{S^{i+j}}{\chi(S)-t}
=[S^{-1}]\frac{S^{i+j}}{\chi(S)}
\quad(0\le i,j<q)
\tag{RW22}
\]

has determinant \((-1)^{q(q-1)/2}\) and satisfies \(A(t)^T\mathsf S=\mathsf SA(t)\). Its t-independence follows because the difference of the two fractions starts at degree at most \(S^{-2}\). With \(J=R^T\mathsf S\), the computed two-parameter pairing identity is

\[
\boxed{A(t)^T J+J A(\varepsilon t)=kJ.}
\tag{RW23}
\]

Indeed transpose \(RA(t)=(kI-A(\varepsilon t))R\) and apply the residue self-adjointness of \(A(\varepsilon t)\). At t=0 this recovers the additive weight relation. For the actual arithmetic residue with denominator g, one must additionally retain the existing full unit \(j_h(g/h)\) and the actual arithmetic inclusion eta; the polynomial chi residue is not silently substituted for that pairing. No extension of that arithmetic unit as a globally invertible operator on every deformed fibre is claimed.

\subsubsection{6. Position in Deligne's argument and in the current program}\label{endpoint-f1_reflection--position-in-delignes-argument-and-in-the-current-program}

The printed primary text of Weil II, pp.~178, 202, 203 and 206, was read in the supplied S20 source-language records together with the existing primary-scan correction. In (3.3.5), the dual degree and Tate twist give \(-n+(2N-i)-2N=-n-i\). In (3.3.6), opposite bounds then constrain the image of the actual compact-support comparison. Page 203 (3.2.13) tensors the same eigenvalue with itself; its square is not replaced by an unspecified tensor constituent.

Our computation supplies a concrete part of the analogous tau-base diagram: the original moment-weight dual action has been transported through the exact exponential deformation and its reflection, including the full special-fibre blocks. It has not supplied Deligne's geometric upper weight bound for the theta quotient. The canonical arithmetic metrics G\_N and their original relation volumes remain unchanged.

The next use of this result is to transport the original residue-dual map and its arithmetic norm comparison along this explicit flow; the factor and parameter shift are now calculated, so they need not be guessed or replaced by a generic purity assertion. The external-product construction and signed determinant calculation are the continuations on those two sides.

\subsubsection{Provenance and verification scope}\label{endpoint-f1_reflection--provenance-and-verification-scope}

\begin{itemize}
\tightlist
\item
  Full controlling tau-base note as pinned above; no whole-paper Lean claim.
\item
  Existing exponential delivery \nolinkurl{Tau_Deligne_Exponential_Comparison/RESEARCH_NOTE.md}, section 7 (37)--(42), read in full for the constant residue duality; DS66--69 supply the previously constructed family and original theta map.
\item
  The scaling calculation supplies the independent chain-level lift underlying (RW14). Formulas (RW6)--(RW13), (RW16), and their joining to the actual weight line are derived here, with their complete algebra above.
\item
  Exact finite fixtures are supporting checks, not replacements for these general proofs. No new Lean execution, RH proof, or remote publication is claimed by this note.
\end{itemize}
